% CHECKPOINT35: complete formalization of the frozen checkpoint34 all-level proof.
% Mathematical predecessor SHA256 0d3055b6d9d98361abdec87d692577f4339f3ae7e963b85c045a92089dbddd6b.
\statement{A.115}{Proved}{The full modular central tower and a compatible peripheral lift}
Retain the original full-\(\widehat{\mathrm{GT}}\) data of A.114. Define
\(L_0=N_1\), \(L_{j+1}=\overline{[L_j,B]}\),
\(L_\infty=\bigcap_j L_j\), and \(A_j=\pi_2(L_j)\).
For every \(j\ge0\),
\[
 \begin{gathered}
 L_j/L_{j+1}\text{ is pro-two},\qquad
 A_{j+1}=\overline{[A_j,\overline B]},\\
 L_{j+1}=L_j\cap\pi_2^{-1}(A_{j+1}),\qquad H_d\in L_j.
 \end{gathered}
\]
The actual word assignments define compatible modular automorphisms
of every \(B/L_j\) and of \(B/L_\infty\), with the original involution
intertwining. A single full peripheral automorphism of \(B\), sending
\(\delta\) literally to \(\delta^d\), induces this whole quotient system.
No cofinality or full equality with the canonical substitution is asserted.
\begin{proof}
\medskip\noindent\emph{1. Series, actual quotients and the precise conclusion}.

Use all actual full-GT canonical data and notation of that proof. In particular
\[
B=C_2(s)\amalg C_3(r),\quad F=\ker(B\to S_3),\quad
N=\ker q,\quad N_1=\overline{[N,F]},
\]
\[
O=\ker(F\to F^{(2)}),\qquad \overline B=B/O,\qquad
M_{21}=\overline{[M_2,F^{(2)}]}.
\]
Define a new, explicitly B-relative series
\[
L_0=N_1,\qquad L_{j+1}=\overline{[L_j,B]},\qquad
L_\infty=\bigcap_{j\ge0}L_j,
\]
and put \(B_j^\flat=B/L_j\), \(A_j=\pi_2(L_j)\le\overline B\). Thus \(B_0^\flat=B_1\) in the notation of A.114; the indexing is not the F-relative series of A.107.

For every finite j we prove:
\[
A_0=M_{21},\qquad A_{j+1}=\overline{[A_j,\overline B]},       \tag{A.115.1}
\]
\[
L_j/L_{j+1}\text{ is pro-two},\qquad
L_{j+1}=L_j\cap\pi_2^{-1}(A_{j+1}),\qquad H_d\in L_j.       \tag{A.115.2}
\]
Each \(B_j^\flat\) has a canonically identified closed normal subgroup \(D_o\), isomorphic to the odd fiber product in A.114, and
\[
B_j^\flat/D_o=\overline B/A_j.                              \tag{A.115.3}
\]
The word assignments \(s\mapsto v=u^{-1}s\), \(\delta\mapsto\delta^d\) give genuine automorphisms \(\Psi_j\) of all these quotients, and of \(B/L_\infty\). They satisfy the actual involution intertwining there.

There exists a full peripheral automorphism \(\widetilde\Psi_\infty\) of B, with cusp \(\delta\mapsto\delta^d\), inducing these same compatible \(\Psi_j\) and hence their limit. This is an existence theorem for a full lift of the quotient actions, not an identification of its full pure action with the canonical \(\psi_d\), nor a full conjugacy of the original \(\alpha\) with \(c_B\).

In particular no cofinality assertion is made. Already
\[
N_1\cap O\subseteq L_\infty                                  \tag{A.115.4}
\]
follows from (A.115.2), so the construction explicitly retains that possible residual.

\medskip\noindent\emph{2. The odd group has unique splittings of central pro-two extensions}.

The actual \(D_o\) from A.114 has
\[
H^1(D_o,\mathbb Q_2/\mathbb Z_2)
=H^2(D_o,\mathbb Q_2/\mathbb Z_2)=0.
\]
The short exact sequence with multiplication by \(2^n\) on this divisible module gives
\[
H^1(D_o,\mathbb Z/2^n)=H^2(D_o,\mathbb Z/2^n)=0.             \tag{A.115.5}
\]
For \(H^1\), use surjectivity of multiplication by \(2^n\) on the invariant trivial module and the zero \(H^1\) above. For \(H^2\), use the two adjacent zero \(H^1\)/\(H^2\) terms. Finite direct sums give (A.115.5) for every finite abelian two-group with trivial action.

Consequently every central extension of \(D_o\) by such a finite group has a unique continuous splitting. Existence is the zero continuous \(H^2\) class; uniqueness is the zero \(H^1\), equivalently the absence of homomorphisms from \(D_o\) to the finite kernel. A finite-kernel profinite extension has a continuous set section: choose an open subgroup meeting its finite kernel trivially, then use finitely many clopen cosets. Thus the usual continuous-cohomology classification applies.

The same assertion holds for a compact abelian pro-two central kernel Z. Quotient the extension by every open subgroup U of Z. Centrality makes U normal in the extension, and the resulting kernel Z/U is finite. Its unique splitting is compatible with every further quotient by uniqueness. Taking the inverse limit gives a continuous splitting of the original extension, because the extension is the inverse limit of these compact quotients and their kernels intersect in one. Uniqueness also follows from (A.115.5), by passing to every finite Z/U.

If this extension is the preimage of a normal copy of \(D_o\) in a larger compact group, its split copy is normal in that larger group. Conjugation transports the splitting, after correcting the base automorphism, to another splitting of the same extension. Uniqueness makes the copies equal. No arbitrary choice of incompatible finite complements is made.

\medskip\noindent\emph{3. Induction of the canonical odd subgroup and odd \(H^2\) vanishing}.

The base \(B_0^\flat=B_1\) has the exact normal \(D_o\) and quotient \(\overline B/M_{21}\) by A.114. It also has zero odd-primary \(H^2\).

Suppose these assertions hold at stage j. Every B-invariant odd-primary character of \(L_j\) transgresses into
\[
H^2(B_j^\flat,\mathbb Q_\ell/\mathbb Z_\ell)=0
\quad(\ell\text{ odd}).
\]
Five-term exactness therefore makes it a restriction of a character of B. It vanishes since \(L_j\subseteq N_1\subseteq B'\). Continuous character separation shows that
\[
Z_j=L_j/L_{j+1}
\]
is pro-two. It is central in \(B_{j+1}^\flat\) by definition.

Restrict the extension \(B_{j+1}^\flat\to B_j^\flat\) to the canonical \(D_o\) at stage j. Subpart 2 gives a unique split copy \(D_o\) at stage j+1, and its normality. The new copy maps isomorphically to the old one. The quotient of the pure subgroup \(F/L_{j+1}\) by this copy is a central pro-two extension of the previous pure pro-two quotient. It is therefore pro-two.

This identifies the odd subgroup with the actual image of O. Indeed the map from F to that pro-two quotient kills O by the defining universal property of the maximal pro-two quotient. Thus the image of O at stage j+1 lies in the new \(D_o\). Its projection to the old \(D_o\) is onto by induction, while the new \(D_o\) maps isomorphically to the old. Hence the image of O is exactly the new \(D_o\). We obtain the literal quotient identity
\[
B_{j+1}^\flat/D_o
 =B/(L_{j+1}O)=\overline B/\pi_2(L_{j+1}).                   \tag{A.115.6}
\]
This step is why the lower comparison continues to be the same actual \(\overline B\), rather than a newly chosen pro-two extension.

Finally the central kernel \(Z_j\) has zero positive cohomology with odd-primary discrete coefficients, by finite-two-group averaging and passage to continuous limits. Inflation therefore preserves the vanishing
\[
H^2(B_{j+1}^\flat,\mathbb Q_\ell/\mathbb Z_\ell)=0
\quad(\ell\text{ odd}).
\]
This completes the induction of the group structure and the odd vanishing.

\medskip\noindent\emph{4. Every next coordinate kernel and the same actual hexagon}.

The whole proof of (A.114.7)--(A.114.8) in A.114 applies to the actual \(D_o\) at every stage:
\[
H^2(\overline B/A_j,\mathbb Q_2/\mathbb Z_2)
 \xrightarrow{\ \mathrm{inf}\ }
H^2(B_j^\flat,\mathbb Q_2/\mathbb Z_2)
\quad\text{is an isomorphism}.                              \tag{A.115.7}
\]
We retain the already proved
\(H^2(\overline B,\mathbb Q_2/\mathbb Z_2)=0\).
For the actual extension diagram
\[
\begin{array}{ccccccccc}
1&\to&L_j&\to&B&\to&B_j^\flat&\to&1\\
 &&\downarrow&&\downarrow&&\downarrow\\
1&\to&A_j&\to&\overline B&\to&\overline B/A_j&\to&1
\end{array}
\]
the same five-term naturality shows that every B-invariant two-primary character of \(L_j\) is
\[
\chi=\chi_{2,j}\circ\pi_2|_{L_j}+\lambda|_{L_j},              \tag{A.115.8}
\]
where \(\chi_{2,j}\) is \(\overline B\)-invariant on \(A_j\) and \(\lambda\) is a character of B. The odd characters were already shown to vanish.

Continuous images of closed commutator subgroups under a compact surjection give
\[
\pi_2(L_{j+1})
 =\overline{[\pi_2(L_j),\overline B]},
\]
proving (A.115.1). Commutator functoriality gives the forward inclusion in (A.115.2). For its reverse, (A.115.8) and the vanishing odd characters kill every element of
\(L_j\cap\pi_2^{-1}([A_j,\overline B])\); character separation of \(L_j/L_{j+1}\) gives exactly \(L_{j+1}\). In particular
\[
L_j/L_{j+1}\simeq A_j/A_{j+1}
\]
as actual compact groups.

A.106 gives \(H_d\in L_0\), and A.89.3 gives its entire actual pro-two image one for the same canonical u,d. The exact coordinate equality now proves inductively \(H_d\in L_j\) for every j. It also proves (A.115.4) immediately. Nothing in this argument makes that intersection trivial.

\medskip\noindent\emph{5. Genuine automorphisms, factor frames and all finite-stage full lifts}.

Because \(H_d\in L_j\), the exact ordered cubic identity gives a homomorphism
\[
B\to B_j^\flat,\qquad s\mapsto v,\quad r\mapsto v\delta^d.
\]
Its pure restriction is the actual \(\psi_d\). This homomorphism kills \(L_j\). To see this without assuming a full source endomorphism of B, start with \(\psi_d(N_1)\subseteq N_1\), already proved in A.107. The homomorphism sends each successive commutator subgroup \([L_i,B]\) into the corresponding image of \(L_{i+1}\). At i=j it therefore kills \(L_j\). It descends to \(\Psi_j\).

If \(N^{F}_0=N\), \(N^{F}_{i+1}=[N^{F}_i,F]\) is the earlier F-relative series, then
\[
N^{F}_{j+1}\subseteq L_j.
\]
This follows from equality at j=0 and commutator induction. A.107's automorphy on \(F/N^{F}_{j+1}\) makes the induced pure map on \(F/L_j\) surjective; it is therefore an automorphism by finite generation and Hopficity. The map on \(S_3\) is the identity. Thus \(\Psi_j\) is an automorphism of \(B_j^\flat\). The same actual generator identity gives \(\alpha\Psi_j=\Psi_jc_B\). The quotient maps commute with all these automorphisms.

For every finite j, the kernel of \(B_j^\flat\to B_1\) is pro-two: it is a finite iterated extension of the central pro-two groups in (A.115.2). Lift the A.108 r-frame from \(B_1\) to stage j. The residual order-three cocycle is in this actual pro-two kernel. The complete coprime \(C_3\) lemma of A.108 kills it, providing a factor frame upstairs. Equivalently apply the same correction at each central step. The s-frame exists already in the full source from v\textsuperscript{2}=1 and the finite-factor torsion theorem. The cusp is literal \(\delta^d\).

In particular stage j+1 is the positive-peripheral central refinement required for lifting stage j, after the same actual multiplier reduction as in A.114. Use an actual multiplier-d Galois automorphism. The original LM7\_PUSH E65.15 covariance makes it preserve N, \(N_1\) and hence every \(L_j\); the original cusp formula is literal. Composition with its inverse removes the common unit power, including the r-branch. Apply the closed-normal modular central criterion of A.3/A.34 to the kernel \(L_j\) and its next closed commutator \(\overline{[L_j,B]}=L_{j+1}\). It gives a full peripheral automorphism inducing \(\Psi_j\). Its cusp conjugator can be corrected inside \(L_j\) so that the cusp is literally \(\delta^d\), as in A.114.

No unrestricted order-preserving or equivariant lifting theorem is used. The conclusion is a lift of each specified quotient automorphism, with its peripheral data.

\medskip\noindent\emph{6. Compact compatibility of the full lifts}.

The actual quotient B/\(L_\infty\) is the inverse limit of B/\(L_j\): compatible cosets have a common point by compactness. The compatible automorphisms \(\Psi_j\) and their inverses therefore give an automorphism \(\Psi_\infty\) of this actual quotient.

The group Aut(B) is compact. One concrete justification is the standard bounded-order construction of cofinal fully invariant open subgroups in a finitely generated profinite group: only finitely many homomorphisms to finite groups of bounded order occur. The endomorphism generator-image space is compact, and invertibility on all these finite quotients is a closed condition. The peripheral conditions, the fixed multiplier d and the literal cusp equation are also closed, since the required conjugators range over compact B.

For each j, let \(C_j\) be the closed set of full peripheral automorphisms inducing \(\Psi_j\) and sending \(\delta\) literally to \(\delta^d\). Subpart 5 proves \(C_j\) nonempty. The sets are descending, because the quotient automorphisms are compatible. Compactness gives
\[
\widetilde\Psi_\infty\in\bigcap_j C_j.
\]
It induces \(\Psi_\infty\) and stabilizes \(L_\infty\). This proves the asserted compatible full lift, without a prior choice of mutually compatible finite-stage conjugators.

Its conjugate \(\widetilde\Psi_\infty c_B\widetilde\Psi_\infty^{-1}\) agrees with the original \(\alpha\) on B/\(L_\infty\). Equality on B has not been shown. Similarly its pure restriction agrees with the canonical \(\psi_d\) only on the displayed quotient system. The residual subgroup need not be cofinal or trivial, and the original coherent-frame/CM obstruction remains undetermined.

\end{proof}
\noindent\emph{Dependencies and scope.}
The proof uses the actual signed marking of A.103--A.108,
the complete A.3/A.34 closed-normal modular lifting criterion,
and the original \(\mathrm{LM7\_PUSH}\) E.65.8/E.65.15 cusp and covariance
formulas. Its continuous cohomological input is Beyl's Proposition~5.7(a,b),
with the full product and lattice calculation printed above or in A.114.
The original full involution and coherent frame problems remain open.
